\section{Implementation and validation}
\label{sec:impl}

\paragraph{Commander4.} In our branch of Commander4 the spectral-index Metropolis step of the
per-pixel path takes a \texttt{likelihood} option. \texttt{marginal} evaluates
Eq.~(\ref{eq:Lflat}) and \texttt{ridge} Eq.~(\ref{eq:Lprof}) in the acceptance, with the
amplitudes integrated or maximized analytically per pixel in a compiled kernel; \texttt{conditional}
is the Gibbs step at fixed amplitudes; \texttt{resolve\_chisq} keeps the move of
Sect.~\ref{sec:legacy} with a warning; \texttt{auto} resolves to \texttt{ridge} where the kernel
applies and to \texttt{conditional} otherwise. Each exact option is tested against a grid evaluation
of its target in unit tests. End to end, a seven-rank run on simulated maps matched the exact profile
posterior within 0.03 posterior standard deviations with width ratios 0.94 and 0.97
(\exref{E006}); in that realization the target itself sat 1.9 and 3.7 standard deviations from the
input indices, which a 40-seed repetition showed to be scatter of the harness at its noise level and
not an offset (\exref{E037}). On real nside-64 intensity maps the \texttt{ridge} chain matched an independent grid
evaluation of its target with shifts of $-0.06 \pm 0.04$ and $+0.06 \pm 0.05$ standard deviations and
width ratios 0.96 and 0.94 for the two indices (\exref{E014}). On the same maps the conditional Gibbs
step has a standard deviation 180 times smaller than the marginal one for the synchrotron index and
did not move in 2000 iterations, while the collapsed chains mix in tens of iterations: collapsing the
amplitudes is needed for mixing as well as for the target. \pend{Simulation-based calibration of the
ridge and marginal modes in Commander4 itself (300 replicates) is not yet run.}

\paragraph{The low-resolution module.} The analysis of Sect.~\ref{sec:data} uses a separate, small
dense implementation (\texttt{nextgen/lowres}): band-averaged spectra from the instrument
bandpasses, the collapsed objectives of Sect.~\ref{sec:model} for several coupled regions with dense
noise covariances, analytic gradients, and $\J$ and $\Tm$ of Eq.~(\ref{eq:JT}). At
$N_{\rm side} = 16$ with three components the amplitude system has 18\,432 unknowns, and a profile fit
of four regions takes three to five minutes on a desktop machine. Unit tests
check the gradients against finite differences, the diagonal-noise closed forms, $\J$ and $\Tm$
against $\sum u V$ and $\sum V$, the sandwich against the scatter of 300 fits, the plug-in bias and
its correction, and the identity of Eq.~(\ref{eq:tilt}) to nine digits. The bandpass readers
reproduce the published LFI colour-correction ratios to $3\times10^{-4}$ and the HFI unit conversions
to 0.3 percent (\exref{E038}), and the per-radiometer effective frequencies of the WMAP team's own
routine to 0.005 GHz (\exref{E039}). An independent statistical review of the module, by a reviewer who derived the targets from the
stated model before reading our derivations, found no code defect; it corrected the response formula
of Sect.~\ref{sec:flat} (the curvature of the marginal in place of the profile information), the
condition under which the profile score is unbiased when a component carries a prior
(App.~\ref{app:sandwich}), and a docstring, and it measured the limits of the plug-in error quoted in
Sect.~\ref{sec:profile} (\exref{E039}, independent review).
